% s-241-p241-m002; original PDF 241
\Aphorism{s-241-p241-m002}{334}
Wer viele Farben auf seiner Palette, viele Formen und Situationen im Stifte hat, arbeitet für die träge Beschauerphantasie zu rasch. Zu gewissen Arten des Stiles, die ein Prestissimo im Geniessen voraussetzen, bringen nur Wenige die geeigneten und geschulten Organe mit; der Normalmensch muss in ein Bild oder eine Stimmung hineinpräludirt, lange darin festgehalten, langsam wieder hinausbegleitet werden, wenn man auf ihn wirken will. Das Schauen einer flüchtig sich abspielenden Bilderfolge, von der gleichwohl jedes ein einheitlich empfundenes, rein ausgearbeitetes, mit den Vor- und Rückbeziehungen reichlich durchsetztes Genrestück darstellt, verlangt schnelles, folgsames, und doch im Augenblicke sich ganz vertiefendes Miterleben. Die meisten Klagen über unruhigen, flackernden, bilderwirrigen Stil entspringen daraus, dass das mangelhaft disciplinirte Ohr dem \Foreign{tempo} nicht folgen kann, weit seltener daraus, dass der Künstler selbst sich nicht folgen kann, dass er die raschen Fiorituren verwischt, durch überstürzte Modulation die Harmonie verdunkelt, die Bilder und Stimmungen, die er auf einer Perlenschnur anreihen sollte, zu einem Haufen durcheinanderschüttet.
\par

